\section{Monitoraggio del traffico di rete}

La raccolta di dati da fonti eterogenee è oggi un'operazione relativamente semplice;
la vera sfida risiede nell'interpretare ciò che viene raccolto, in particolare
nel dominio della sicurezza informatica, dove le minacce reali rappresentano
eventi statisticamente rari in un flusso di traffico altrimenti monotono.

Poiché la stragrande maggioranza del traffico è benigna e gli attacchi più pericolosi
sono rari e difficili da distinguere dal rumore di fondo, è necessario un framework
strutturato che integri raccolta, archiviazione e analisi in un sistema coeso.

\subsection{Architettura generale di un framework di analisi di rete}

Un framework di analisi di rete ha lo scopo di integrare dati provenienti da fonti eterogenee
per fornire una visione completa e contestualizzata degli eventi che si verificano
all'interno di un'infrastruttura.
Esso orchestra il flusso di informazioni dai sensori di rete fino alla console dell'analista,
abilitando sia analisi in tempo reale che indagini storiche approfondite.

L'architettura generale di tale framework si basa su diversi componenti interconnessi:

\begin{itemize}
    \item \textbf{Sensor Network} (rete di sensori): è il punto di partenza,
    responsabile della raccolta dei dati grezzi della rete, degli host e dei servizi.

    \item I dati possono essere inviati a un sistema di \textit{realtime processing}
    (elaborazione in tempo reale) per analisi immediate,
    come il rilevamento di frame anomali o la generazione di log.

    \item Contemporaneamente, i dati confluiscono nel \textbf{repository},
    il cuore del sistema di archiviazione, che a sua volta si suddivide in:

    \begin{itemize}
        \item \textbf{Archive}: deposito per i dati degli eventi.
        \item \textbf{Annotation}: contenitore per i dati di arricchimento.
        \item \textbf{Knowledge Base} (KB): base di conoscenza specifica dell'organizzazione.
    \end{itemize}

    \item Il sistema di \textbf{Query Processing} (elaborazione delle interrogazioni)
    attinge al repository per permettere agli analisti di effettuare ricerche complesse
    e sintetizzare le informazioni.

    \item Infine, la \textbf{console} funge da interfaccia di interazione utente,
    permettendo di visualizzare i risultati delle analisi in tempo reale e delle interrogazioni,
    e di interagire con il sistema.

\end{itemize}

\subsubsection{Archive}

L'archive è il deposito dedicato ai dati degli eventi.
Questi possono includere una vasta gamma di informazioni, come:

\begin{itemize}
    \item \textbf{Dati di traffico di rete}: log di NetFlow, catture complete dei pacchetti (pcap).
    \item \textbf{Informazioni di sistema}: statistiche sui processi, sull'utilizzo del file system,
    report di antivirus.
    \item \textbf{Log di sistema}: log dei server web, informazioni di syslog.
\end{itemize}

Una caratteristica chiave dell'archive è che i suoi dati sono \textbf{immutabili}:
una volta scritti, non vengono aggiornati.
Questo lo rende l'ambiente ideale per eseguire analisi forensi e storiche,
garantendo la coerenza e l'integrità delle evidenze.
Per ottimizzare lo spazio, per i dati più datati è spesso conveniente conservare
solo i risultati post-elaborazione, anziché i dati grezzi completi.

\subsubsection{Annotation}
